\section{What Would Refute the QR Architecture?}\label{sec:failure}
A proposed local transducer fails when its maps or domains are incompatible, when its required receipt cannot be represented, or when its complete map is many-to-one. A claimed recovery fails concretely if distinct admitted inputs have identical complete outputs. A particular failed construction refutes its dependent claims, not every event architecture.

A finite correspondence fails if its asserted matrix algebra, normalization, or full outcome table is wrong; registration fails if it changes the intended measure or violates its setting interface. Mode fails on a proposed domain if the claimed action is not defined there, its conjugacy fails, or its parity is not a well-defined presentation character. The six-frame result in Section~\ref{sec:mode} is checked on its stated domain.

A physical bridge is rejected when its fixed mapping contradicts observations in its claimed regime. The common-foundation claim is weakened if the two reductions require unrelated adjustable structures. Every bridge must therefore specify its native inputs, target observable, map or limit, and failure criterion: do not protect the idea; make the idea protect itself.

A proposed complexity functional fails its declared presentation invariance if equivalent registered objects receive different values. A proposed local scalar relay fails when identical available data require different increments or target scalars. Missing an initial reference prevents absolute reconstruction from increments alone. These are failures of the specified functional or interface, not counterexamples to the telescoping identity once its premises hold.


A proposed conservative representation fails its stated metric law when its residual $U^\dagger\mathsf G U-\mathsf G$ is nonzero. A proposed Maxwell-type complex fails constraint preservation when its declared divergence-of-curl composition is nonzero. Conversely, a vanishing metric residual does not establish native admission or physical correspondence; the source-selection and representation gates remain independently testable in Section~\ref{sec:nativedynamics}.

A proposed incidence reconciliation fails when the lifted orientation characters disagree on a shared isotropy element, or when the actual common-domain maps fail to intertwine. A changed sign convention cannot repair that failure. A claim that phase transport induces the original history-state exchange fails to be established until the corresponding decorated readout and its intertwiner are supplied; missing evidence is not itself a counterexample to the proposed extension.

\Needspace{0.22\textheight}
\section{Methodological Guardrails}

\label{sec:guardrails}
The following distinctions specify the mathematical interfaces used in this manuscript.

\begin{longtable}{@{}p{0.35\textwidth}p{0.59\textwidth}@{}}
\toprule
Distinction & Consequence \\
\midrule
\endhead
Carrier versus transducer & A graph supplies relations, not a unique event grammar, instrument, or physical theory. \\
Visible versus complete return & A projected endpoint can recur while a lift, receipt, or retained word differs. \\
Event depth versus tick & Primitive handoff cost need not equal oriented winding; a tick can contain several events. \\
Measure versus frequency versus provenance & Correct weights and a realizing sequence do not identify a native chronological source. \\
Receipt versus copied payload & Quantum admissibility forbids arbitrary informative records of an otherwise undisturbed unknown state. \\
Joint information versus marginal sums & Correlations are part of the conserved joint structure; Shannon mutual information is subtracted in the entropy identity. \\
Local controller versus complete machine & Unbounded accessible history makes the complete system more than a closed finite-state automaton. \\
Deck multiplicity versus branching entropy & Constant covering multiplicity does not by itself create an extensive future-growth rate. \\
Carrier symmetry versus execution symmetry & The event grammar may break a symmetry of the unadorned carrier. \\
Closure versus reset & Comparison can leave a residual distinction even when visible closure holds. \\
Spectral sheet versus physical surface & No spatial dimension follows merely from a decomposition into eigenspaces or cover sheets. \\
Common invariant versus full unification & One bridge must be followed by compatible dynamics, observables, and controlled reductions. \\
Native normalization versus calibration & $\cth=1$ does not supply metres, seconds, energy, or a value of $\hbar$. \\
Mode versus phase versus time & Mode is the registered $S_3$ orientation class, phase is $C_3$ position, and $\tth$ is integer cycle winding. Their crosswalks require compatible domains. \\
Mode versus collective eigenmode & Modal parity grades a registered presentation; an eigenmode belongs to a specified dynamical operator. \\
Orientation state versus reversal action & $b/ab$ are torsor members, $Z$ is a grading, and $\Mevt$ is an operation; equal cardinality does not identify them. \\
Common cover versus complete event & A many-to-one registration needs a retained distinction for full recovery; local graph bijectivity does not supply event admission. \\
Family $Q_{30}$ versus original $G_{30}$ & Equal order does not identify quotient maps or transfer the certified temporal visibility theorem. \\
Kernel points versus relation bins & The twelve direction/character bins contain $96$ ordered pairs; they are not the twelve points or twelve Born histories. \\
Frame isotropy versus free deck action & The rooted $D_8$ and the composite deck group $C_2^3$ have different actions despite equal order. \\
Census quotient choice versus QR Mode & Distinct central involution quotients are not automatically two presentations in one modal orbit. \\

Relational complexity versus information & $\mathcal C$ measures registered relational organization; complete information preservation constrains the update, not every scalar readout. \\
Scalar versus phase/Mode & Descriptive orthogonality is not an inner product or independence theorem. Presentation changes preserve scalar complexity only on the declared equivalent-object domain. \\
Scalar relay versus independent payload & Reconstruction needs sufficient available relational data and, for incremental relay, an initial value. Global conservation alone is not a local decoder. \\
Return depth versus general complexity & Linear repeated-return complexity is conditional on the chosen history domain and additive functional; signed charge and distinct support are not interchangeable. \\
Causal unit versus complexity unit & $\cth=1$ does not imply $\Delta\mathcal C=1$. Each normalization requires its own domain and witness. \\
Naming versus evidence & Thales supplies intellectual background and Tolkien a secondary resonance; neither supplies a mathematical or physical premise. \\

Recovery versus quadratic invariance & Neither implies the other for arbitrary maps. A faithful linear conservative flow supplies both only under its declared hypotheses. \\
Native complexity versus physical magnitude & $\|\mathbf L\|$, electromagnetic energy, and $I_{\mathsf G}$ are not identified with $\mathcal C$ by analogy. \\
Tangent motion versus global rotation & A unit-axis constraint forces tangent velocity, not a constant generator, a fixed cone, or a linear isometry. \\
Axis stabilizer versus normal kernel & $SO(3)/SO(2)$ is a homogeneous space; it is not a quotient group like a character quotient. \\
Oriented axis versus unoriented line & Their stabilizers are $SO(2)$ and embedded $O(2)$ respectively. The latter's reversal changes the oriented axis. \\
Fiber phase versus native event phase & Mechanical axial phase needs a full frame and local convention. No identification with native $C_3$ phase or capitalized Mode is made. \\
Skew form versus selected dynamics & A conserved metric does not select a generator or force. Group-averaged finite metrics do not establish native provenance. \\
Maxwell energy versus information current & Field energy, charge current, Poynting flux, and encoded-decodable information transport retain different definitions. \\
Closed field versus subsystem balance & Sources and boundaries contribute explicit work and flux. Enlarging the system requires a model, not merely a conservation slogan. \\
Formal skewness versus skew-adjointness & The field operator's domain and boundary realization are part of the claim; the periodic construction is specified explicitly. \\
Reciprocal exchange versus phase reversal & WXYZTI exchange, modal conjugation, and signed electromagnetic duality have different action laws unless a common-domain map proves otherwise. \\
Finite step versus finite-state dynamics & Matrix exponentiation or interpolation need not remain inside the native admitted action family; a continuum bridge requires an additional construction. \\

Selected reference versus privileged frame & A presentation marking fixes coordinates; it does not derive or select a physical preferred frame. \\
Reference choice versus native coupling & Naming a torsor member positive cannot construct a missing common incidence or repair unequal isotropy characters. \\
Incidence sign versus absolute orientation & Signs compare participants at a specified incidence. Their transformation laws must include both endpoint reference changes. \\
One global bit versus several orbits & Two comparison solutions follow only on a common transitive domain with compatible isotropy. Multiple orbits can carry independent signs. \\
Equal 120 counts versus common objects & Cover vertices, bare incidences, sector permutations and signed operators have different domains and projections. \\
Projective section versus representation & A global factor set can obstruct multiplication of the selected representatives even when their doubled extension splits. \\
Phase actor versus history state & Conjugating word generators does not itself prove a readout intertwiner on the original $\{2,3\}$ states. \\
Passive sign change versus active effect & $(A,a)\mapsto(-A,-a)$ preserves a projector. Reversing an observable on one leg with labels fixed requires a separate admitted action. \\
\bottomrule
\end{longtable}

An implementation can satisfy one row while failing another. For example, a finite table may be exactly quantum and no-signaling while its chronological generator lacks an admitted local implementation. A complete audit therefore follows the interfaces rather than grading a whole narrative with one pass/fail label.
